%========================================================================
\section{Data and Target Construction}
\label{sec:data}

% Flughöhe: Eine kompetente Person soll dies OHNE deinen Code reimplementieren
% können. Beschreibe WAS und WARUM, nicht Funktionssignaturen.

\subsection{Country universe and sample period}
\label{sec:data:universe}

We study monthly sovereign credit spreads for ten euro-area issuers
(Austria, Belgium, Spain, Finland, France, Greece, Ireland, Italy, the
Netherlands, and Portugal). The spread is the difference between the ECB-10y yield of a country and Germany (ECB dataset IRS, series key M.{cc}.L.L40.CI.0000.EUR.N.Z). The sample runs from 2003-04
to 2026-04: the start is the first month, and the end the last month, in which
all required data points are simultaneously available.

Greece is an exception. In 2015-07 there was no market data because of the
capital controls during the Greek crisis. This missing window is carried forward and doesn't lie in the test set. A second limitation of including Greece is the 2012 private-sector
involvement, which produced a structural break in the spread data.

\subsection{Target variable: delta-spreads}
\label{sec:data:target}

Let $\spread_{i,t}$ denote the level credit spread of country $i$ at time
$t$. All models are trained on the \emph{first difference}
\begin{equation}
  \dspread_{i,t} \;=\; \spread_{i,t} - \spread_{i,t-1},
  \label{eq:delta-target}
\end{equation}
and every forecast is mapped back to level space before evaluation via
\begin{equation}
  \hat{\spread}_{i,t} \;=\; \spread_{i,t-1} + \widehat{\dspread}_{i,t}.
  \label{eq:back-transform}
\end{equation}
Spreads are highly persistent, and taking the first difference removes this
dominant persistence and avoids spurious level persistence. Evaluation is still
performed in level space, via \cref{eq:back-transform}.

\subsection{Feature groups and the leakage constraint}
\label{sec:data:features}

The feature selection is not part of this work, so we adhere to the features
proposed in \textcite[ECB WP 1781]{afonso2015} with some additions. This is a
good fit, since \textcite[ECB WP 1781]{afonso2015} investigates spread drivers
for the same countries and the same monthly data structure that we forecast.
One caveat is that they use data from 1999-01 to 2010-12, a period that does not
even cover half of ours, and, as \textcite[ECB WP 1781]{afonso2015} themselves
show, the predictors of credit spreads change over time.

We use monthly industrial production growth as a monthly GDP proxy, since GDP
growth is only available quarterly. This follows \textcite[ECB WP 1781]{afonso2015}.
They use expected debt/GDP and fiscal balance/GDP instead of the realized values
we use. Expected data would probably yield more signal, but they were not
available to us. Because realized data are subject to revisions that inject
information unknown at the reference date, we implement a lag dictionary to
prevent this data leakage (\cref{sec:data:lags}). We would also have liked to
use the bid-ask spread, but here too no free data were available.
\Cref{tab:features_overview} lists all the features we use.

\begin{table}[htbp]
    \centering
    \small
    \caption{Feature catalogue: identifier, construction, scope and primary
    source. Country fundamentals enter as differentials vs.\ Germany (DE);
    global features are identical across all countries.}
    \label{tab:features_overview}
    \resizebox{\textwidth}{!}{%
    \begin{tabular}{@{}clllcl@{}}
        \toprule
        \textbf{\#} & \textbf{Identifier} & \textbf{Description} & \textbf{Construction} & \textbf{Scope} & \textbf{Source} \\
        \midrule
        1  & \texttt{spread\_lag}              & lagged own spread        & $s_{i,t-1}$                                                                  & local  & \parencite{afonso2015} \\
        2  & \texttt{vix\_log}                 & log VIX                  & $\log(\mathrm{VIX})$                                                         & global & \parencite{afonso2015} \\
        3  & \texttt{us\_baa\_aaa}             & Baa$-$Aaa corp.\ spread  & Moody's Baa $-$ Aaa                                                          & global & \parencite{bernoth2004sovereign} \\
        4  & \texttt{us\_aaa\_treasury}        & US corp$-$Treasury       & Moody's Aaa $-$ 10y UST                                                      & global & \parencite{codogno2003yield} \\
        5  & \texttt{epu\_log}                 & log EPU Europe           & $\log(\mathrm{EPU})$                                                         & global & \parencite{baker2016epu} \\
        6  & \texttt{debt\_gdp}                & rel.\ debt/GDP           & $(\text{debt/GDP})_i-(\text{debt/GDP})_{\text{DE}}$                          & local  & \parencite{afonso2015} \\
        7  & \texttt{debt\_gdp\_sq}            & rel.\ debt/GDP squared   & $(\text{debt/GDP}_i-\text{debt/GDP}_{\text{DE}})^{2}$                        & local  & \parencite{afonso2015} \\
        8  & \texttt{vix\_log\_x\_debt}        & intl.\ risk $\times$ debt & $(\texttt{vix\_log}-\text{exp.\ mean})\times\texttt{debt\_gdp}(t{-}4)$        & local  & \parencite{afonso2015} \\
        9  & \texttt{us\_aaa\_treasury\_x\_debt} & US risk $\times$ debt   & $(\texttt{us\_aaa\_treasury}-\text{exp.\ mean})\times\texttt{debt\_gdp}(t{-}4)$ & local  & \parencite{codogno2003yield} \\
        10 & \texttt{fiscal\_balance}          & rel.\ fiscal balance     & 4Q-mean(B9), rel.\ DE                                                        & local  & \parencite{afonso2015} \\
        11 & \texttt{debt\_service\_rev}       & interest / revenue       & 4Q-mean(D41PAY)\,/\,4Q-mean(TR), rel.\ DE                                    & local  & \parencite{bernoth2004sovereign} \\
        12 & \texttt{ip\_growth\_yoy}          & IP growth YoY            & YoY \% growth of IP, rel.\ DE                                                & local  & \parencite{afonso2015} \\
        13 & \texttt{reer\_log}                & log REER                 & $\log(\text{REER})_i-\log(\text{REER})_{\text{DE}}$                          & local  & \parencite{afonso2015} \\
        14 & \texttt{liquidity\_share}         & rel.\ market size        & $\text{debt}_i/\sum_{j}\text{debt}_j$ (11 incl.\ DE)                         & local  & \parencite{bernoth2004sovereign} \\
        15 & \texttt{target2\_gdp}             & TARGET2/GDP              & $(\text{TARGET2}/\text{GDP})\times100$                                       & local  & \parencite{degrauweji2013} \\
        16 & \texttt{pc2\_core\_periph}        & 2nd principal comp.      & PC2 score (expanding, see text)                                             & global & \parencite{afonso2015} \\
        17 & \texttt{rating\_num}              & sovereign rating         & $\text{rating}_i-\text{rating}_{\text{DE}}$ (1--17)                          & local  & \parencite{afonso2012sovereign} \\
        \bottomrule
    \end{tabular}%
    }
\end{table}

\paragraph{Global risk factors.}
\textcite[ECB WP 1781]{afonso2015} uses only the VIX as a global risk factor;
we add three more. \textcite{codogno2003yield} showed that US credit risk is a
spread driver in its own right, so we include \texttt{us\_aaa\_treasury}
(Moody's Seasoned Aaa Corporate Yield $-$ 10y US Treasury), as well as the
deviation of \texttt{us\_aaa\_treasury} from its average times the difference in
debt/GDP ratio to Germany, as proposed by \textcite{codogno2003yield}; entering
the deviation rather than the level avoids collinearity of this term with the
debt term. \textcite{bernoth2004sovereign} identified Moody's Baa Corporate
Yield $-$ Moody's Aaa Corporate Yield as an indicator of investor risk aversion
and, by that, a driver of sovereign credit spreads. We also include the Economic
Policy Uncertainty Europe index as a log. It measures how much newspapers report
on economic uncertainty; its relevance in this context was shown by
\textcite{bernal2016epu}.
\footnote{We aimed to include the $us\_hy\_oas$ = US High Yield Option-Adjusted Spread as an indicator of the global risk tolerance, but without a license it was only available three years into the past.}

\paragraph{Fiscal fundamentals and relative market size.}
\textcite{bernoth2004sovereign} also identified interest expenditure / total
government revenue (as a difference to Germany) and the outstanding debt of the
country / sum of the debt of the 11 panel countries (incl.\ DE) as sovereign
credit spread drivers, so we include them as well. The fiscal balance and the
debt-service ratio have a strong seasonal character, as seen from a regression
on seasonal dummies with $R^{2}=0.35$. We therefore replace the quarterly data
by their trailing four-quarter mean, which removes the seasonal component
(quarter-dummy $R^{2}\approx0$) while keeping the four-month publication lag. This had to be done since the Seasonally and Calendar Adjusted data was not available here and we had to resort to Non-Seasonally-Adjusted data. The industrial production (IP) index is available as a SCA series, while the fiscal block (government balance B9, interest payable D41PAY, and total revenue TR) relies on non-seasonally adjusted (NSA) data, since Eurostat (e.g., table $gov\_10q\_ggnfa$) does not provide adjusted series for all countries; we apply the trailing four-quarter moving average there as well. Other variables like the real effective exchange rate (REER) and sovereign debt levels inherently do not require seasonal adjustment. For the debt-service ratio, the numerator and denominator each receive the same
treatment before dividing. The Eurostat quarterly government finance statistics
are forward-filled until the next datapoint is available.

\paragraph{Sovereign rating.}
\textcite{afonso2012sovereign} has shown that the numeric sovereign rating on a
scale of 1--17, as a difference to Germany
($\text{rating}_i - \text{rating}_{\text{DE}}$), is also relevant, so we include
it.

\paragraph{Convertibility and capital flight.}
The Trans-European Automated Real-time Gross settlement Express Transfer System,
2nd generation, short TARGET2, is a market-based indicator of capital flight.
\textcite{degrauweji2013} showed that the loss of investor confidence led to a
bond sell-off and then a liquidity outflow in 2010--11. TARGET2 is what measures
exactly the accounting counterpart of this liquidity outflow, as shown by
\textcite{sinnwollmershaeuser2012,krishnamurthy2018}. We use
$(\text{TARGET2}/\text{GDP})\times100$. \textcite{degrauweji2013} do not use
TARGET2 as an empirical proxy for their theory. GDP itself is reconstructed as
$\texttt{gov\_debt\_mio\_eur} / (\texttt{gov\_debt\_pct\_gdp}/100)$, since we
already have both series.

\paragraph{Country-specific risk.}
The global risk factors above are identical across countries. To capture
country-specific risk, we also implement a feature used in \textcite{afonso2015}:
the deviation of the log VIX from its expanding mean, times the relative debt per
country. Relative debt is again debt/GDP minus the same for Germany, and we again
work with the realized debt instead of the expected.

\paragraph{Cross-sectional heterogeneity.}
\texttt{pc2\_core\_periph} is calculated from the z-standardised country data,
with a 12-month refit on an expanding window of minimum 36 months. The loadings
are anchored so that periphery countries always load positive. The missing value
for Greece is handled by imputing 0. PC2 measures only the core--periphery
contrast and is not a measure of monetary risk.

Overall, the global features are the same for every country, while the local
features change from country to country. Each country's data contains all the
global features as well as the features specific to that country, so every
country contributes the same 17 features.

\subsection{Publication-lag protocol}
\label{sec:data:lags}

As already stated, we did not have access to vintage data and had to resort to later published, revised data instead. In the revision process, information can enter that was not yet available in the reference period. That is a form of data leakage we want to prevent. So we employ a lag dictionary: every datapoint in reference period $t$ is moved $k=\lceil \text{lag}/30\rceil$ full months forward, so we can only use the datapoint at $t + k$. Since we are always snapping to the end of the month, we risk overlagging by 29 to 30 days. The target itself is never lagged. Look-ahead bias still remains in the monthly average series (Moody's Aaa/Baa yields, EPU): they appear about 5 days into the next month. We still employ no time lag since we judge
it negligible and record it here rather than leave it unstated.


% Präambel: \usepackage{booktabs}
\begin{table}[H]
\centering
\caption{Publication lag by feature. The shift we impose in real time is
$k=\lceil \text{days}/30\rceil$ whole months on a month-end series
(see Note).}
\label{tab:publication-lags}
\begin{tabular}{@{}p{4.3cm}p{1.2cm}p{1.4cm}p{5.3cm}@{}}
\toprule
\textbf{Feature(s)} & \textbf{Lag (d)} & \textbf{Shift (m)} & \textbf{Source} \\
\midrule
\texttt{vix\_log}                                   & 0    & 0  & \parencite{fred_vixcls} \\
\texttt{epu\_log}$^{\dagger}$                        & 0    & 0  & \parencite[p.\,6, fn.\,8]{baker2016epu} \\
\texttt{spread\_lag}, \texttt{pc2\_core\_periph}     & 0    & 0  & \parencite{ecb_irs_calendar} \\
\texttt{us\_aaa\_treasury}$^{\dagger}$, \texttt{us\_baa\_aaa}$^{\dagger}$
                                                     & 0    & 0  & \parencite{fed_h15_qa} \\
\texttt{rating\_num}                                 & 0    & 0  & \parencite{spglobal_emea_calendar,afonsofurcerigomes2012} \\
\texttt{reer\_log}                                   & 30   & 1  & \parencite{eurostat_erteff} \\
\texttt{target2\_gdp}$^{\ddagger}$                    & 33   & 2  & \parencite{ecb_target_balances} \\
\texttt{ip\_growth\_yoy}                             & 45   & 2  & \parencite{eurostat_ip_releases} \\
\texttt{debt\_gdp}, \texttt{debt\_gdp\_sq},
\texttt{liquidity\_share}                            & 113  & 4  & \parencite{eurostat_gov10qggdebt} \\
\texttt{fiscal\_balance}, \texttt{debt\_service\_rev}& 113  & 4  & \parencite{eurostat_gov10qggnfa} \\
\texttt{vix\_log\_x\_debt},
\texttt{us\_aaa\_treasury\_x\_debt}$^{\ddagger}$      & --   & -- & \parencite{eurostat_gov10qggdebt} \\
\bottomrule
\end{tabular}
\vspace{4pt}
\raggedright\footnotesize
\textit{Note:} \emph{Lag (d)} reports the provider's documented publication
delay after the reference period. \emph{Shift (m)} is the lag we actually
impose, $k=\lceil \text{days}/30\rceil$ whole months, applied as a month-end
shift of the feature within each country. The forecast target is never
shifted. Because $k$ rounds up, no value enters the panel before it was
released; the price we pay is up to roughly 29 days of over-lagging. \\[2pt]
$^{\dagger}$ The monthly-average series (Moody's Aaa/Baa yields, EPU) appear
about 5 days into the following month. At monthly frequency that delay rounds
to a zero-month shift, so a look-ahead of a few days remains. We judge it
negligible and record it here rather than leave it unstated. \\[2pt]
$^{\ddagger}$ \emph{Split lag.} When an input combines two series with
different release calendars, we lag each series on its own calendar before
combining them. For \texttt{target2\_gdp} $=$ TARGET2$/$GDP the TARGET2
numerator is shifted by 2 months and the GDP denominator (from quarterly EDP
debt) by 4 months, which gives $\text{TARGET2}(t{-}2)/\text{GDP}(t{-}4)$. This
holds the numerator, the part that moves fastest in a crisis, at its shortest
defensible lag. The interaction terms follow the same logic: in
$\text{global}(t)\times\text{debt}(t{-}4)$ the debt factor already carries its
4-month (113 d) lag before it multiplies the contemporaneous global factor
(lag 0), so we do not lag the product a second time.
\end{table}

\subsection{Final panel}
\label{sec:data:panel}
The final panel covers all 10 countries over 277 months, from 2003-04 to
2026-04, giving a balanced $10 \times 277 = 2770$ country-month observations.
Each country carries the same 17 features, of which 5 are global (identical
across countries) and 12 are local (country-specific). The panel is stored in
wide format: the index is the month-end date, and the columns are a MultiIndex
of country and variable. It carries the spread level \emph{unshifted} and holds
no pre-shifted target; the forecasting label $\dspread_{i,t+1}$ is constructed
model-side (\cref{sec:data:target}). The only missing value is the Greek spread
in 2015-07. The temporal train/test structure and the sequence construction are
described in \cref{sec:model-setup}.


